\subsection{张量代数的正向极限}

设 \((\mathbf{A}_{\alpha},\phi_{\beta \alpha})\) 为一个交换环的有向正向系统，且 \((\mathbf{M}_{\alpha},f_{\beta \alpha})\) 为一个由 \(\mathbf{A}_{\alpha}\) -模构成的正向系统（II，§ 6，第 2 号）；设 \(\mathbf{A} = \varinjlim \mathbf{A}_{\alpha}\) 和 \(\mathbf{M} = \varinjlim \mathbf{M}_{\alpha}\) ，后者是一个 A-模。对于 \(\alpha \leqslant \beta\) ，一个 \(\mathbf{A}_{\alpha}\) -代数同态（第 3 号，公式 (8)） \(f_{\beta \alpha}' = \mathsf{T}(f_{\beta \alpha})\colon \mathsf{T}_{\mathbf{A}_{\alpha}}(\mathbf{M}_{\alpha})\to \mathsf{T}_{\mathbf{A}_{\beta}}(\mathbf{M}_{\beta})\) 由 \(\mathbf{A}_{\alpha}\) -同态 \(f_{\beta \alpha}\colon \mathbf{M}_{\alpha}\to \mathbf{M}_{\beta}\) 典范地导出，并且由（9）（第 3 号）可知 \((\mathsf{T}_{\mathbf{A}_{\alpha}}(\mathbf{M}_{\alpha}),f_{\beta \alpha}')\) 是一个由 \(\mathbf{A}_{\alpha}\) -代数构成的正向系统。另一方面，设 \(f_{\alpha}\colon \mathbf{M}_{\alpha}\to \mathbf{M}\) 为典范 \(\mathbf{A}_{\alpha}\) -同态；一个 \(\mathbf{A}_{\alpha}\) -代数同态 \(f_{\alpha}'\colon \mathsf{T}_{\mathbf{A}_{\alpha}}(\mathbf{M}_{\alpha})\to \mathsf{T}_{\mathbf{A}}(\mathbf{M})\) 由（第 3 号，公式 (8)）导出，并且由（9）（第 3 号）也可得出，诸 \(f_{\alpha}'\) 构成一个由 \(\mathbf{A}_{\alpha}\) -同态构成的正向系统。

命题6. A-同态 \(f' = \varinjlim f_{\alpha}' : \varinjlim T_{A_{\alpha}}(M_{\alpha}) \to T_{A}(M)\) 是一个分次代数同构。

为简化起见，我们记 \(\mathbf{E} = T_{\mathbf{A}}(\mathbf{M})\) 和 \(\mathbf{E}' = \varinjlim T_{\mathbf{A}_{\alpha}}(\mathbf{M}_{\alpha})\) ，并设

\[
g _ {\alpha}: \mathsf {T} _ {\mathrm{A} _ {\alpha}} (\mathbf {M} _ {\alpha}) \rightarrow \mathbf {E} ^ {\prime}
\]

为典范 \(\mathbf{A}_{\alpha}\) -同态。显然，复合 \(\mathbf{A}_{\alpha}\) -线性映射 \(\mathbf{M}_{\alpha} \xrightarrow{\phi_{\mathbf{M}_{\alpha}}} \mathsf{T}_{\mathbf{A}_{\alpha}}(\mathbf{M}_{\alpha}) \xrightarrow{g_{\alpha}} \mathbf{E}'\) 构成一个正向系统，因此存在唯一的 A-线性映射 \(u = \varinjlim (g_{\alpha} \circ \phi_{\mathbf{M}_{\alpha}}): \mathbf{M} \to \mathbf{E}'\) ，使得

\[
u \circ f _ {\alpha} = g _ {\alpha} \circ \phi_ {M _ {\alpha}}
\]

对所有 \(\alpha\) 成立。该映射本身唯一地分解（第 1 号，命题 1）为 \(\mathbf{M} \xrightarrow{\phi_{\mathbf{M}}} \mathrm{E} \to \mathrm{E}'\) ，其中 \(h\) 是一个 A-代数同态。只需证明 \(h \circ f' = 1_{\mathrm{E}'}\) 和 \(f' \circ h = 1_{\mathrm{E}}\) .

为此注意到，对于每个指标 \(\alpha\) （第3号，公式（8））

\[
h \circ f _ {\alpha} ^ {\prime} \circ \phi_ {M _ {\alpha}} = h \circ \phi_ {M} \circ f _ {\alpha} = u \circ f _ {\alpha} = g _ {\alpha} \circ \phi_ {M _ {\alpha}}
\]

由此，根据第1号命题1的唯一性断言，

\[
h \circ f _ {\alpha} ^ {\prime} = g _ {\alpha}
\]

对所有 \(\alpha\) 成立；由此可得 \((h \circ f') \circ g_{\alpha} = g_{\alpha}\) 对所有 \(\alpha\) 成立，因此由正向极限的定义得 \(h \circ f' = 1_{\mathbf{E}'}\) 。

另一方面，根据第3号公式（8），

\[
f ^ {\prime} \circ u \circ f _ {\alpha} = f ^ {\prime} \circ g _ {\alpha} \circ \phi_ {M _ {\alpha}} = f _ {\alpha} ^ {\prime} \circ \phi_ {M _ {\alpha}} = \phi_ {M} \circ f _ {\alpha},
\]

由此再次得到 \(f' \circ u = \phi_{\mathbb{M}}\) （由正向极限的定义）；我们得出 \(f' \circ h \circ \phi_{\mathbb{M}} = \phi_{\mathbb{M}}\) ，而第 1 号命题 1 的唯一性性质给出 \(f' \circ h = 1_{\mathbb{E}}\) .

命题6也可以通过观察以下事实来证明：对于每个整数 \(n \geqslant 1\) ，存在一个典范的 A-模同构 \(\lim_{\longrightarrow} \mathsf{T}_{\mathbf{A}_{\alpha}}^{n}(\mathbf{M}_{\alpha}) \to \mathsf{T}_{\mathbf{A}}^{n}(\mathbf{M})\) ，这由 II，§6，第3号，命题7通过对 \(n\) 作归纳得到。立即可验证，这些同构都是 \(f'\) 的限制。

